% Last Updated: March 1, 2024
\documentclass[0-main]{subfiles}

\begin{document}
\ifSubfilesClassLoaded{
	% \title{\tbf{Untitled}}
	% \author{Cong Wen, xx/20xx}
	% \maketitle
	\tableofcontents
}{}
% ----------------------------------------------------------------------------------------------------


\ssc{Inverse Galois problem}
Ref:~\cite{MM2018}

\sss{Rigidity method}

Inverse Galois problem asks which finite groups are realizable as Galois groups over a given field $K$, the case of most interest is when $K=\bQ$. \tbf{Rigidity method}, which is the emphasis of the book, aims at solving the middle arrow of the following descents, since the first and last arrow are known by theorems of Grothendieck and Hilbert.
\[
	\begin{tikzcd}
		\bC(t)\ar[r, dash] & \oL{\bQ}(t)\ar[r, dash] & \bQ(t)\ar[r, dash] & \bQ
	\end{tikzcd}
\]


\prg{The case $\bC(t)$: Motivation}
\begin{thm}[Theorem~1.3, Profinite Riemann Existence Theorem]
	Let $X=\bP^{1}$, $S\sbs X(\bC)$ a finite set, denote by $\bC(t)^{S}$ the maximal Galois extension of $\bC(t)$ unramified away from $S$. Then one has
	\[
		\begin{tikzcd}
			\pi_{1}^{\mathrm{top}}(X(\bC)\ssm S)\ar[r, hook, "\ita"] & \oH{\pi}_{1}^{\mathrm{top}}(X(\bC)\ssm S)\ar[r, "\cong"] & \pi_{1}^{\mathrm{alg}}(X(\bC)\ssm S)=\Gal(\bC(t)^{S}/\bC(t)).
		\end{tikzcd}
	\]
\end{thm}

\begin{cor}
	Any finite group can be realized as a Galois group over $\bC(t)$.
\end{cor}

\prg{From $\bC(t)$ to $\oL{\bQ}(t)$: Theorem of Grothendieck}
\begin{thm}[{\cite[XIII, Corollaire~2.12]{SGA1}, Theorem~2.2}]
	Let $X=\bP^{1}$, $k\sbs\bC$ an algebraically closed field, let $S\sbs X(k)$ be a finite set, denote by $k(t)^{S}$ the maximal Galois extension of $k(t)$ unramified away from $S$. Then one has
	\[
		\begin{tikzcd}
			\pi_{1}^{\mathrm{top}}(X(\bC)\ssm S)\ar[r, hook, "\ita"] & \oH{\pi}_{1}^{\mathrm{top}}(X(\bC)\ssm S)\ar[r, "\cong"] & \pi_{1}^{\mathrm{alg}}(X(k)\ssm S)=\Gal(k(t)^{S}/k(t)).
		\end{tikzcd}
	\]
\end{thm}

\begin{cor}
	Any finite group can be realized as a Galois group over $\oL{\bQ}(t)$.
\end{cor}

\prg{From $\bQ(t)$ to $\bQ$: Hilbert irreducibility}
\begin{thm}[{\cite[Theorem~\Rnum{1}]{Hilbert1892}}]
	If $f(x,t)\in\bQ[x,t]$ is irreducible, then there are infinite many choices of $a\in\bQ$ such that $f(x,a)\in\bQ[x]$ is irreducible.
\end{thm}

\begin{cor}
	Any finite group realizable over $\bQ(t)$ is realizable over $\bQ$.
\end{cor}

\prg{From $\oL{\bQ}(t)$ to $\bQ(t)$: Rigidity method}

\begin{thm}[Theorem~2.4, Theorem~2.6, Splitting Theorem]
	Let $X=\bP^{1}$, let $S\sbs X(\oL{\bQ})$ be a finite set invariant under $\Gal(\oL{\bQ}/\bQ)$, then $\oL{\bQ}(t)^{S}$ is Galois over $\bQ(t)$, and
	\[
		\Gal(\oL{\bQ}(t)^{S}/\bQ(t))\cong\Gal(\oL{\bQ}(t)^{S}/\oL{\bQ}(t))\rtm\Gal(\oL{\bQ}/\bQ).
	\]
	Furthermore, the action of the semidirect product is given by
	\[
		[\gma_{i}^{\oL{\dta}}]=[\gma_{\dta(i)}^{\chi(\dta)}], \dta\in\Gal(\oL{\bQ}/\bQ).
	\]
	where $\oL{\dta}$ is a lift of $\dta$, $\chi$ the cyclotomic character, equality is as conjugacy classes.
\end{thm}

\begin{rmk}
	To see the splitting, it suffices to note $\Gal(\oL{\bQ}/\bQ)$ is isomorphic to the decomposition group of a degree $1$ prime divisor $\kP$ whose extension $\oH{\kP}$ to $\oL{\bQ}(t)^{S}$ is unramified.

	Easier example is the case $\bR(t)$, where one only has to lift the complex conjugation $\sA{\sgm}$, to describe its action on the topological generators, if they are preserved by $\sA{\sgm}$, then those loops in pair are send to each other, and those single loops is send to its inverse conjugated by other single loops before it.
\end{rmk}

\begin{thm}[{\cite[Page~71]{Serre1997}, Belyi's Theorem\cite{Belyi1979}}]
	Let $X$ be a smooth projective curve over $\bC$, then the following are equivalent,
	\begin{enr}
		\item $X$ can be defined over $\oL{\bQ}$.
		\item There is a finite index subgroup $\Gma\lqs\SL_{2}(\bZ)$ such that $X\cong(\bH\cup\bP^{1}(\bQ))/\Gma$.
		\item There is a finite cover $X\ar\bP^{1}(\bC)$ unramified away from $\{0,1,\ift\}$.
	\end{enr}
\end{thm}

\begin{thm}[Theorem~4.1, Hurwitz Classification]
	Let $X=\bP^{1}$, $S\sbs X(\oL{\bQ})$ a finite set with $\sP{S}=s$, consider
	\begin{align*}
	  \Sgm_{S}(G) &= \{\sgm=(\sgm_{1},\ldots,\sgm_{s})\in G^{s}\mid \sA{\sgm_{1},\ldots,\sgm_{s}}=G, \sgm_{1}\cdots\sgm_{s}=1\} \\
	  N_{S}(G) &= \{N\mid \oL{\bQ}(t)^{S}\gqs N\gqs \oL{\bQ}(t), \Gal(N/\oL{\bQ}(t))\cong G\}
	\end{align*}
	then $\Sgm_{S}(G)$ admits an action of $\Aut(G)$, and there is a surjection
	\begin{align*}
		\Sgm_{S}(G)/\Inn(G)&\ar N_{S}(G) \\
		\sgm=(\sgm_{1},\ldots,\sgm_{s})&\amt N_{\sgm}=(\oL{\bQ}(t)^{S})^{\Ker\psi_{\sgm}}
	\end{align*}
	where $\psi_{\sgm}\in\Hom{}{}{(\Gma_{s},G)}$ is the map sending $\gma_{i}$ to $\sgm_{i}$.
\end{thm}

\begin{thm}[Proposition~4.4]
	Let $G$ be a finite group, and $C\in \mathrm{Cl}(G)^{s}$. One can define
	\[
		\Dta_{C} = \{\dta\in\Gal(\oL{\bQ}/\bQ)\mid C^{\dta}=C\}, \bQ_{C} = \oL{\bQ}^{\Dta_{C}}=\bQ(\{\chi(C)\mid\chi\in\Irr(G)\}).
	\]
	Then $\Dta_{C}\sbs\Gal(\oL{\bQ}/\bQ)$ is closed, and $\bQ_{C}\sbs\bQ(\zta_{\sP{G}})$ is abelian, with index and degree $d(C)$, where
	\[
		d(C)=\sP{\{C^{n}\mid n\in (\bZ/\sP{G}\bZ)\xT\}}.
	\]
	Say $C$ is rational if $d(C)=1$.
\end{thm}

\begin{thm}[Theorem~4.5]
	Let $G$ be a finite group, and $C\in \mathrm{Cl}(G)^{s}$, moreover take $\sgm\in\Sgm_{S}(G)\cap C$ such that the ramification locus of $N_{\sgm}/\oL{\bQ}(t)$ is defined over $\bQ$. One can define
	\[
	  \Dta_{\sgm} = \{\dta\in\Gal(\oL{\bQ}/\bQ)\mid \sgm^{\dta}=\sgm\}, K_{\sgm} = \oL{\bQ}(t)^{\Dta_{\sgm}}.
	 \]
	Then $[K_{\sgm}:\bQ_{C}(t)]\leq l(C)$, where
	\[
		l(C)=\sP{(\Sgm_{S}(G)\cap C)/\Inn(G)}.
	\]
	Say $C$ is rigid if $l(C)=1$.
\end{thm}

\begin{thm}[Proposition~4.9]
	Let $G$ be a finite group, and $C\in\Cl(G)^{s}$, one can further define
	\[
		\Sym(C)=\{\oga\in S_{s}\mid \xis n\in(\bZ/\sP{G}\bZ)\xT, C^{\oga}=C^{n}\}.
	\]
	Then any subgroup $V\lqs\Sym(C)$ defines $C^{V}=\{C^{\oga}\mid \oga\in V\}$, and
	\[
		\Dta_{C}^{V}=\{\dta\in\Gal(\oL{\bQ}/\bQ)\mid C^{\dta}\in C^{V}\}, \bQ_{C}^{V}=\oL{\bQ}^{\Dta_{C}^{V}}.
	\]
	Then $\Dta_{C}^{V}\sbs\Gal(\oL{\bQ}/\bQ)$ is closed, $\bQ_{C}^{V}\sbs\bQ(\zta_{\sP{G}})$ is abelian, with index and degree $d^{V}(C)=d(C)/\sP{C^{V}}$, where
	\[
		d^{V}(C)=d(C)/\sP{C^{V}}.
	\]
	Say $C$ is $V$-symmetric $d^{V}(C)=1$.
\end{thm}

\begin{rmk}
	In practice, given a finite simple group $G$ and class vector $C\in\Cl(G)^{s}$, one would hope to check rationality and rigidity effectively.

	For rationality, one has to understand all irreducible characters. On can use \cite{CCNPW1985} for sporadic simple groups, and use Deligne-Lusztig theory for finite groups of Lie type \cite{GM2020}.

	For rigidity, verifying $l(C)=1$ is clear from definition, there are also many criterions concerning as below.
\end{rmk}

\begin{thm}[Theorem~4.8, Basic Rigidity]
	Let $G$ be a finite group whose center has a complement, and there exists a rational rigid class vector $C\in\Cl(G)^{s}$. Then $G$ can be realized over $\bQ(t)$.
\end{thm}

\begin{thm}[Theorem~4.11, Strong Rigidity]
	Let $G$ be a finite group whose center has a complement, and there exists a rigid class vector $C\in \mathrm{Cl}(G)^{s}$. Let $V$ be a symmetry group of $C$ such that for any $\dta\in\Dta_{C}^{V}$ there exists a unique $\oga\in V$ with $C^{c(\dta)}=C^{\oga}$. Then $G$ can be realized over $\bC_{C}^{V}(t)$.
\end{thm}

\prg{Examples concerning rigidity and rationality}


\begin{thm}[Theorem~5.1]
	Every finite abelian group can be realized over $\bQ$.
\end{thm}

\begin{prf}[Sketch of the proof]
	It suffices to deal with the cyclic group $C_{n}=\sA{\sgm}$.
	\begin{itm}
		\item If $n=2$, $(\sgm,\sgm\xI)$ is rigid and rational.
		\item If $n\gqs3$, $(\sgm^{k}\mid k\in(\bZ/n\bZ)\xT)$ is rigid and can be made $V$-symmetric by a suitable choice of $V$.
	\end{itm}
	So by Basic Rigidity and Strong Rigidity, they can all be realized over $\bQ$.
\end{prf}


\begin{thm}[Proposition~5.2]
	$S_{n},A_{n},n\gqs3$ can be realized over $\bQ$.
\end{thm}
\begin{prf}[Sketch of the proof]
	It suffices to find rational and rigid class vector in them.

	For $S_{n}$, $C=(2A, (n-1)A, nA)$ is rigid, since it must be conjugate to $((12),(2,\ldots,n)\xI,(1,\ldots,n))$, and every class vector of $S_{n}$ is rational, hence by Basic Rigidity.

	For $A_{n}$, consider the realized finite Galois extension $N/\bQ(t)$ with Galois group $S_{n}$. Then $K=N^{A_{n}}$ is ramified over $\bQ(t)$ at only two points, by Riemann-Hurwitz, $K$ is a rational function field over $\bQ$, i.e., $K\cong\bQ(s)$, hence $A_{n}$ is realizable.
\end{prf}


\prg{Rigidity Criterion}

\begin{thm}[Corollary~5.9]
	For a finite group $G$, consider a class vector $C$ with some $\sgm\in C\cap\Sgm_{S}(G)\neq\ept$, $C$ is rigid if
	\[
		\sum_{\chi\in\Irr(G)}\frac{\chi(\sgm_{1})\cdots\chi(\sgm_{s})}{\chi(1)^{s-2}}=\frac{\sP{C_{G}(\sgm_{1})}\cdots\sP{C_{G}(\sgm_{s})}}{\sP{G}^{s-2}\sP{Z(G)}}.
	\]
\end{thm}

\begin{rmk}
	This applies to $\SL_{2}(\bF_{8})$, by using ATLAS\cite{CCNPW1985}.
\end{rmk}


\begin{thm}[Theorem~5.10, Rigidity Criterion]
	Let $G$ be a finite group, $k$ any field, $R:G\ar\GL_{n}(k)$ an irreducible faithful representation. Suppose
	\begin{enr}
		\item there exists $\sgm_{1},\sgm_{2}\in G$ such that $\sA{\sgm_{1},\sgm_{2}}=G$,
		\item there exists $a\in k\xT$ such that $\sgm_{1}-a\Id$ has rank $1$,
		\item $N_{\GL_{n}(k)}(G)=G\cdot C_{\GL_{n}(k)}(G)$
	\end{enr}
	Then any $([\sgm_{1}],[\sgm_{2}],[\sgm_{2}\xI\sgm_{1}\xI])$ is a rigid triple.
\end{thm}

\sss{Guralnick-Malle and Yun}
Guralnick and Malle proved the inverse Galois problem for $E_{8}(p), p\geq7$ using rigidity method\cite{GM2014}, inspired by Yun's motive paper\cite{Yun2014Motives}.


\sss{Non-normal extensions}
Ref:~\cite{FK1978}

Note that the inverse Galois problem over $\bQ$ asks which finite groups are realizable as the automorphism group of a \tbf{Galois extension} of $\bQ$. \tbf{Fried-Koll{\'a}r} solved this if the Galois condition, in fact, the normal condition, is dropped in an extremely short paper\cite{FK1978}. We collect the construction here.
% \begin{lem}\label{FK1978:Lem1}
% 	Let $p$ be a rational prime, $L$ a field, $a\in L\bs L^{p}$. If some $b\in L$ satisfies $b^{1/p}\in L(a^{1/p})$, then there exists $c\in L$ and $0\leq k<p$, $b^{1/p}=ca^{k/p}$.
% \end{lem}

% \begin{prf}
% 	Note that $b^{1/p}=f(a^{1/p})$ for some $f(x)\in L[x], \Deg f<p$. Now $L(\zta_{p},a^{1/p})$ is Galois over $L(\zta_{p})$, $f(\zta_{p}a^{1/p})=\zta_{p}^{k}b^{1/p}$ for some $0\leq k<p$, hence $a^{1/p}$ is a root of $g(x)=f(\zta_{p}x)-\zta_{p}^{k}f(x)\in L(\zta_{p})[x]$, note $x^{p}-a$ is irreducible over $L(\zta_{p})$, so $g(x)=0$, hence the result.
% \end{prf}

\begin{lem}\label{FK1978:Lem2}
	Let $L$ be a number field, $\cO_{L}$ its ring of integers. Consider any finite set of monic polynomials $\{f_{i}\}_{1\leq i\leq n}\sbs\cO_{L}[x]$ and all of them are not $p$-th powers in $\cO_{L}[x]$ for some prime $p$, then there exists $t\in \cO_{L}$ such that all $\{f_{i}(t)\}_{1\leq i\leq n}$ are not $p$-th powers in $\cO_{L}$.
\end{lem}

\begin{prf}
	The smooth projective curves $V(y^{p}-f_{i}(x))\sbs\bP_{L}^{2}$ all have genus $g>0$ by assumption. Siegel's theorem says they have finitely many integral points. $\cO_{L}$ is infinite, so the desired $t$ exists.
\end{prf}

\begin{thm}
	For any finite group $G$, there exists a finite extension $K_{G}/\bQ$ such that $\Aut(K_{G}/\bQ)\cong G$.
\end{thm}

\begin{prf}[Sketch of the proof]
	$K_{G}$ is constructed in $5$ steps:
	\begin{enr}
		\item $G\cong \Aut(X)$, where $X=(V,E)$ is an acyclic undirected connected graph, assume $\#V=n\gqs5$.
		\item $S_{n}\cong\Gal(L/\bQ)$, where $L$ is the splitting field of $F(x)$ with algebraic integers $\{a_{i}\}_{1\leq i\leq n}$ as roots.
		\item For a prime $p\geq5$, take $t\in\cO_{L}$ for the polynomials $\prod_{1\leq i<j\leq n}(x+a_{i}+a_{i})^{k_{ij}}, 0\leq k_{ij}<p$ by Lemma~\ref{FK1978:Lem2}.
		\item Choose a root $b_{ij}$ for each polynomial $x^{p}-(a_{i}+a_{j}+t), 1\leq i<j\leq n$.
		\item $K_{G}=L(\{b_{ij}\}_{(i,j)\in E})$ is the desired field with $\Aut(K_{G}/\bQ)\cong G$.
	\end{enr}
	Note that the additional $k_{ij}$-s are used to guarantee the last step, more precisely, to show $\zta_{p}b_{k}\nin L(b_{1},\ldots,b_{k-1})$.
	% There are two key points to see this:
	% \begin{enr}
	% 	\item Use Lemma~\ref{FK1978:Lem1} to show $\zta_{p}b_{k}\nin L(b_{1},\ldots,b_{k-1})$ for each $k$.
	% 	\item None of the fields $L(b_{1},\ldots,b_{k})$ contains $\zta_{p}$, hence they are not normal over $\bQ$.
	% \end{enr}

	% Then verifying $\Aut(K_{G}/\bQ)\cong G$ is direct.
\end{prf}







% ----------------------------------------------------------------------------------------------------
\ifSubfilesClassLoaded{
	\bibliographystyle{amsalpha}
	\bibliography{/Users/wenc/Documents/Math/universal-ref}
}{}
\end{document}
